% Draft manuscript for the A2/R-157/R-158 publication lane.
% This file is intentionally a manual synthesis of the registered proof notes.
\documentclass[11pt]{article}
\usepackage[a4paper,margin=1in]{geometry}
\usepackage{amsmath,amssymb,amsthm,mathtools}
\usepackage{booktabs}
\usepackage{enumitem}
\usepackage[T1]{fontenc}
\usepackage{lmodern}
\usepackage{hyperref}

\hypersetup{colorlinks=true,linkcolor=blue,citecolor=blue,urlcolor=blue}
\numberwithin{equation}{section}
\newtheorem{theorem}{Theorem}[section]
\newtheorem{proposition}[theorem]{Proposition}
\newtheorem{lemma}[theorem]{Lemma}
\newtheorem{corollary}[theorem]{Corollary}
\theoremstyle{definition}
\newtheorem{definition}[theorem]{Definition}
\newtheorem{remark}[theorem]{Remark}
\newcommand{\T}{\mathbb T}
\newcommand{\C}{\mathbb C}
\newcommand{\R}{\mathbb R}
\newcommand{\cF}{\mathcal F}
\newcommand{\cO}{\mathcal O}
\newcommand{\Id}{\mathrm I}
\newcommand{\dd}{\,\mathrm d}
\newcommand{\norm}[1]{\left\lVert#1\right\rVert}
\newcommand{\ip}[2]{\left\langle #1,#2\right\rangle}
\newcommand{\grad}{\nabla_{\!\R}}
\newcommand{\eps}{\varepsilon}

\title{Global Well-Posedness, Neutral-State Rejection, and
       Ensemble-Induced Shell Coexistence in a Regularized
       Multicomponent Brazovskii Functional}
\author{Jusang Lee}
\date{Draft v0.1.0 -- 2026-09-03}

\begin{document}
\maketitle

\begin{abstract}
We study one explicitly defined three-component complex field functional on a
fixed periodic three-torus.  The field is treated as six real components and
the derivative Class-II terms are regularized by a positive density floor.
Conditional only on the registered identification of this functional with the
canonical P1 production functional, we prove global well-posedness of its
real-$L^2$ gradient flow for every $H^2$ datum, exact energy dissipation,
continuous dependence, and positive-time smoothness.  An exact completion of
the quadratic and local polynomial terms, together with a positive radial
derivative of the Class-II energy, shows that the unconstrained functional has
the zero field as its unique critical point and global minimizer and that the
gradient flow decays exponentially in $L^2$.  We then solve a distinct,
imposed ensemble problem.  On the side-$16$ torus, the exact spectral bottom
lies on the integer shell $|n|^2=3$.  After fixing
$Q=\norm{\Psi}_{2}^{2}/2$ or introducing
$\cO_\mu=\cF-\mu Q$, an exact nonnegative decomposition gives a finite-shell
constrained minimizer and a first-order grand-potential coexistence point.
The charge or chemical potential is imposed rather than derived; the paper
therefore makes no claim about a physical vacuum, BCC structure, infinite
volume, or a quantum continuum limit.  We provide a claim-to-proof map and a
reproducibility protocol tied to the repository's independent exact audits.
\end{abstract}

\section{Introduction and statement of scope}

There are two mathematically different questions in the current TECT/A2
candidate.  The first is whether a fixed regularized functional defines a
well-behaved dissipative evolution and what its unconstrained equilibria are.
The second is what happens after the variational problem is changed by fixing
a nonzero quadratic charge or by subtracting a chemical-potential term.  A
nonzero minimizer of the second problem must not be reported as an equilibrium
of the first problem.  This distinction is the organizing principle of the
present paper.

The object is a specific functional, not an assertion that this functional is
the physical law of nature.  Its identification with the canonical P1
production functional is retained as the named hypothesis
\texttt{A2-H3-CANONICAL-PRODUCTION-FUNCTIONAL}.  All conclusions below are
conditional on that identification and unconditional for the explicitly
defined mathematical functional once the identification is accepted.

The paper has three main contributions.
\begin{enumerate}[label=(\roman*)]
\item We give a continuum proof for the full regularized three-component
      fourth-order functional, including the derivative Class-II terms.  The
      nonlinear map has order two, not four, and the endpoint $H^4$ estimate
      is obtained by a Duhamel cancellation rather than by an invalid
      fractional-smoothing shortcut.
\item We give an exact global sign test for the pinned unconstrained problem.
      It rejects every nonzero equilibrium of this candidate, not merely the
      finite list of lattice patterns tested previously.
\item We solve the separately declared fixed-charge and grand-canonical
      problems by an exact spectral and polynomial completion.  The result is
      a finite-volume ensemble mechanism, with an explicit discontinuous
      charge jump, while the original neutral reference remains lower.
\end{enumerate}

The statements are deliberately narrower than a constructive quantum-field
theory theorem.  We do not claim an infinite-volume limit, a real-time
automorphism group, algebraic KMS dynamics, a ground-state gap, a continuum
renormalization, a derived conserved charge, a density crystal, BCC selection,
or any cosmological interpretation.

\section{The functional and conventions}

Let $\T^3_{16}=[0,16)^3$ with periodic boundary conditions and let
$\Psi:\T^3_{16}\to\C^3$.  We regard $\C^3$ as $\R^6$ and use the real pairing
\begin{equation}
 \ip{U}{V}_{\R}=\operatorname{Re}\int_{\T^3_{16}}U^\dagger V\,\dd x.
 \label{eq:real-pairing}
\end{equation}
Write $\rho=\Psi^\dagger\Psi$.  For the three Hermitian internal generators
$T_A$ in the pinned P1 data, set
\begin{equation}
 m_A=\Psi^\dagger T_A\Psi,\qquad
 J_A=\nabla m_A,\qquad
 K_A=J_A-\frac{m_A}{\rho+\eps_\rho}\nabla\rho,
 \qquad \eps_\rho=10^{-12}.
 \label{eq:currents}
\end{equation}
The positive floor in \eqref{eq:currents} is part of the declared functional;
removing it is a different theorem.

The quadratic operator is
\begin{equation}
 L=\bigl(r+Z(-\Delta)+Y(-\Delta)^2\bigr)\Id_3+M_{\rm fam}
       +k_{\rm lock}(\Id_3-P_0),
 \label{eq:L}
\end{equation}
where the pinned scalar coefficients are
\begin{equation}
 Y=1,\qquad Z=-0.9252754126,\qquad r=0.4740336473,
 \label{eq:scalar-coeff}
\end{equation}
and the family/lock contribution is positive semidefinite.  The local
polynomial coefficients are
\begin{equation}
 \lambda=-\frac{43}{100},\qquad \gamma=\frac{81}{50}>0.
 \label{eq:local-coeff}
\end{equation}
The Class-II density is
\begin{equation}
 \cF_{\rm II}(\Psi)=\sum_A\int_{\T^3_{16}}
 \left(\frac a2|J_A|^2+b\operatorname{Re}(J_A^\dagger K_A)
                  +\frac c2|K_A|^2\right)\dd x,
 \label{eq:classii}
\end{equation}
with the pinned positive-definite coefficient matrix
$Q_{\rm II}=\left(\begin{smallmatrix}a&b\\b&c\end{smallmatrix}\right)$.
The full functional is
\begin{equation}
 \cF(\Psi)=\frac12\ip{\Psi}{L\Psi}_{\R}
 +\int_{\T^3_{16}}\left(\frac\lambda4\rho^2+\frac\gamma6\rho^3\right)\dd x
 +\cF_{\rm II}(\Psi).
 \label{eq:functional}
\end{equation}

The domain for the evolution theorem is $H^2(\T^3_{16};\C^3)$, while the
linear operator has domain $H^4$ and form domain $H^2$.  All assertions about
the charge ensemble below refer to the same fixed torus and the same
functional \eqref{eq:functional}; no new physical field or reservoir is
introduced.

\section{Main results}

\begin{theorem}[full gradient-flow well-posedness]
Assume the named canonical-functional hypothesis H3.  For every
$\Psi_0\in H^2(\T^3_{16};\C^3)$, the real-$L^2$ gradient flow
\begin{equation}
 \partial_t\Psi=-\grad\cF(\Psi),\qquad \Psi(0)=\Psi_0,
 \label{eq:flow}
\end{equation}
has a unique global solution with
\begin{equation}
 \Psi\in C([0,\infty);H^2)\cap L^2_{\rm loc}((0,\infty);H^4),qquad
 \partial_t\Psi\in L^2_{\rm loc}((0,\infty);L^2).
 \label{eq:solution-class}
\end{equation}
The solution depends continuously on the initial datum in $H^2$ on every
finite time interval and satisfies
\begin{equation}
 \cF(\Psi(t))+\int_s^t\norm{\partial_\tau\Psi(\tau)}_2^2\dd\tau
       =\cF(\Psi(s)),\qquad 0\le s\le t.
 \label{eq:energy-id}
\end{equation}
For every $t>0$ it is smooth in space and time.
\end{theorem}

\begin{theorem}[neutral minimizer and exponential decay]
For the pinned coefficients, there are exact rational constants
\begin{equation}
 g=\frac{719818750025582338837}{5400000000000000000000}>\frac18,qquad
 \kappa=\frac{2101675000076747016511}{8100000000000000000000}>\frac14
 \label{eq:g-kappa}
\end{equation}
such that every $\Psi\in H^2$ obeys
\begin{equation}
 \cF(\Psi)\ge g\norm{\Psi}_2^2,qquad
 \ip{D\cF(\Psi)}{\Psi}_{\R}\ge\kappa\norm{\Psi}_2^2.
 \label{eq:neutral-bounds}
\end{equation}
Consequently $\Psi=0$ is the unique critical point and global minimizer of
the unconstrained problem, and every solution of \eqref{eq:flow} obeys
\begin{equation}
 \norm{\Psi(t)}_2^2\le e^{-2\kappa t}\norm{\Psi_0}_2^2.
 \label{eq:decay}
\end{equation}
\end{theorem}

\begin{theorem}[imposed charge and chemical-potential ensemble]
Let $Q(\Psi)=\norm{\Psi}_2^2/2$ and let $\lambda_0$ be the exact bottom of
the quadratic operator \eqref{eq:L} on the side-$16$ torus.  Put
\begin{equation}
 \rho_* = \frac{43}{216},\qquad
 c_* = \frac{1849}{86400},\qquad \mu_t=\lambda_0-c_*.
 \label{eq:ensemble-constants}
\end{equation}
Then, with $\rho=\Psi^\dagger\Psi$,
\begin{equation}
 \cF(\Psi)=\mu_t Q(\Psi)
 +\frac12\ip{\Psi}{(L-\lambda_0)\Psi}_{\R}
 +\cF_{\rm II}(\Psi)
 +\frac\gamma6\int_{\T^3_{16}}\rho(\rho-\rho_*)^2\dd x.
 \label{eq:ensemble-decomp}
\end{equation}
Every term after $\mu_tQ$ is nonnegative.  A ground-shell plane wave with
constant density $\rho_*$ is a constrained global minimizer at
\begin{equation}
 Q_*=\frac{11008}{27}.
 \label{eq:Qstar}
\end{equation}
For $\cO_\mu=\cF-\mu Q$, zero is the unique minimizer for $\mu<\mu_t$,
coexists with nonzero ground-shell minimizers at $\mu=\mu_t$, and is beaten
by a nonzero minimizer for $\mu>\mu_t$.  The coexistence point lies strictly
between the amplitude saddle-node and the zero-field spinodal:
\begin{equation}
 \mu_{\rm sn}=\lambda_0-\frac{1849}{64800}<
 \mu_t=\lambda_0-\frac{1849}{86400}<\lambda_0=\mu_{\rm lin}.
 \label{eq:first-order-ordering}
\end{equation}
This is a theorem about an imposed ensemble, not a derived physical charge or
neutral-vacuum selection.
\end{theorem}

\section{Proof of the evolution theorem}
\label{sec:evolution-proof}

\subsection{Linear operator and coercivity}

The scalar symbol of \eqref{eq:L} is $p(x)=r+Zx+Yx^2$ with $x=|k|^2$.
Completing the square gives
\begin{equation}
 p(x)=Y\left(x+\frac{Z}{2Y}\right)^2+r-\frac{Z^2}{4Y},qquad
 \inf_{x\ge0}p(x)=0.260000000009475>0.
 \label{eq:symbol-positive}
\end{equation}
Moreover,
\begin{equation}
 \inf_{x\ge0}\frac{p(x)}{1+x^2}=0.2048572626782363>0.
 \label{eq:H2-coercivity}
\end{equation}
The family matrix and the lock term are positive semidefinite, so $L$ is
self-adjoint on $H^4$ and its form domain is $H^2$, with equivalent graph and
Sobolev norms.  The pinned Class-II matrix is positive definite; its recorded
determinant and least eigenvalue are respectively
\begin{equation}
 \det Q_{\rm II}=7.031249999996483\times10^{-6},qquad
 \lambda_{\min}(Q_{\rm II})=0.001259011500926061.
 \label{eq:QII-numbers}
\end{equation}

Since $\lambda<0$ and $\gamma>0$, Young's inequality gives
\begin{equation}
 \frac{|\lambda|}{4}\rho^2\le \frac\gamma{12}\rho^3+C_Y,qquad
 C_Y=0.010098435197886498.
 \label{eq:young}
\end{equation}
Combining \eqref{eq:H2-coercivity}, \eqref{eq:QII-numbers}, and
\eqref{eq:young} yields, for a positive $c_{H^2}$,
\begin{equation}
 \cF(\Psi)\ge \frac{c_{H^2}}2\norm{\Psi}_{H^2}^2
       +\frac\gamma{12}\norm{\Psi}_{L^6}^6-C_Y|\T^3_{16}|.
 \label{eq:energy-coercivity}
\end{equation}

\subsection{The Class-II map has order two}

Realify $\Psi$ as $u\in\R^6$.  Let $A_A$ be the real symmetric realification
of $T_A$, and define
\begin{equation}
 p_A(u)=2A_Au,qquad q_A(u)=\frac{u^TA_Au}{u^Tu+\eps_\rho},qquad
 s_A(u)=2(A_A-q_A(u)\Id)u.
 \label{eq:realified-coeff}
\end{equation}
The Class-II density can then be written exactly as
\begin{equation}
 G_{\rm II}(u,\nabla u)=\frac12\sum_j(\partial_j u)^T B(u)(\partial_j u),
 \quad
 B(u)=\sum_A\{a p_Ap_A^T+b(p_As_A^T+s_Ap_A^T)+c s_As_A^T\}.
 \label{eq:B-map}
\end{equation}
The positive $\eps_\rho$ makes $B$ smooth on all of $\R^6$.  Direct
Euler--Lagrange differentiation therefore has the form
\begin{equation}
 N_{\rm II}(u)=B(u)\nabla^2u+C(u)[\nabla u,\nabla u],
 \label{eq:N-map}
\end{equation}
where $C$ is smooth and no derivative above order two occurs.  On the fixed
three-torus,
$H^2\hookrightarrow L^\infty\cap W^{1,6}$ and $L^3\hookrightarrow L^2$.
Product estimates and the algebra property of $H^2$ thus give, on every
bounded $H^2$ ball,
\begin{equation}
 \norm{N(u)-N(v)}_2\le C_R\norm{u-v}_{H^2}.
 \label{eq:local-lipschitz}
\end{equation}

The positive fourth-order operator generates an analytic semigroup.  In
particular,
\begin{equation}
 \norm{L^{1/2}e^{-tL}}_{L^2\to L^2}\le Ct^{-1/2},
 \qquad t>0.
 \label{eq:semigroup}
\end{equation}
The kernel in \eqref{eq:semigroup} is integrable at the origin.  The mild
equation, \eqref{eq:local-lipschitz}, and a contraction argument in
$C([0,T];H^2)$ yield a unique local solution and the usual continuation
alternative.

\subsection{Galerkin passage and energy identity}

Let $P_n$ be the real $L^2$ Fourier projection and solve
\begin{equation}
 \partial_tu_n=-P_n(Lu_n+N(u_n)).
 \label{eq:galerkin}
\end{equation}
This is the exact finite-dimensional real gradient flow, hence
\begin{equation}
 \frac{\dd}{\dd t}\cF(u_n)=-\norm{\partial_tu_n}_2^2.
 \label{eq:galerkin-energy}
\end{equation}
The coercivity estimate gives a uniform $L^\infty(0,T;H^2)$ bound and
\eqref{eq:galerkin-energy} gives an $L^2(0,T;L^2)$ bound on $\partial_tu_n$.
Ellipticity applied to $Lu_n=-\partial_tu_n-P_nN(u_n)$ then gives an
$L^2(0,T;H^4)$ bound.  Aubin--Lions compactness yields strong convergence in
$L^2(0,T;H^2)$, and \eqref{eq:local-lipschitz} gives strong convergence of
the nonlinear term in $L^2(0,T;L^2)$.

For the nonlinear part $\Phi$, the derivative identity
$D\Phi(u)[h]=\ip{N(u)}{h}_{\R}$ and the local Lipschitz estimate give the chain
rule
\begin{equation}
\Phi(u(t))-\Phi(u(s))=\int_s^t\ip{N(u(\tau))}{\partial_\tau u(\tau)}_{\R}\dd\tau.
 \label{eq:nonlinear-chain}
\end{equation}
The corresponding spectral quadratic-form chain rule combined with
\eqref{eq:nonlinear-chain} proves the equality \eqref{eq:energy-id}, not just
an energy inequality.  The lower bound \eqref{eq:energy-coercivity} prevents
the $H^2$ norm from diverging at a finite maximal time, so the solution is
global.

\subsection{Continuous dependence and smoothing}

For two solutions in a common energy-bounded ball, the mild difference equation
has the Volterra estimate
\begin{equation}
 \norm{u(t)-v(t)}_{H^2}\le C\norm{u_0-v_0}_{H^2}
 +C_R\int_0^t(t-s)^{-1/2}\norm{u(s)-v(s)}_{H^2}\dd s.
 \label{eq:singular-gronwall}
\end{equation}
Weakly singular Gronwall proves continuous dependence on finite intervals.
Analytic-semigroup estimates first give $u(t)\in D(L^\alpha)=H^{4\alpha}$ for
$t>0$ and every $\alpha<1$.  Choose $\alpha>1/2$ to obtain Hölder continuity
into $H^2$ away from $t=0$.  Thus $f(t)=N(u(t))$ is Hölder continuous into
$L^2$, and the endpoint cancellation
\begin{equation}
 L\int_a^t e^{-(t-s)L}f(s)\dd s
 =(\Id-e^{-(t-a)L})f(t)
 +\int_a^tLe^{-(t-s)L}[f(s)-f(t)]\dd s
 \label{eq:endpoint-cancellation}
\end{equation}
has an integrable last term.  It follows that $u(t)\in D(L)=H^4$ for every
$t>0$.  Periodic Moser estimates give
\begin{equation}
 N:H^m\longrightarrow H^{m-2}\quad\text{locally Lipschitz},qquad m\ge2,
 \label{eq:moser}
\end{equation}
and iteration of \eqref{eq:endpoint-cancellation} proves the claimed
$C^\infty$ regularity.

\section{Exact neutral-state result}
\label{sec:neutral-proof}

\subsection{Quadratic completion}

The scalar symbol has the exact minimum
\begin{equation}
 \mu_{\rm sh}=r-\frac{Z^2}{4Y}
 =\frac{26000000000947494031}{10^{20}}>0.
 \label{eq:mu-shell}
\end{equation}
The pinned internal matrix is
\begin{equation}
 M=\begin{pmatrix}
 1/10&-1/20&-1/20\\
 -1/20&13/100&-1/20\\
 -1/20&-1/20&17/100
 \end{pmatrix}.
 \label{eq:M}
\end{equation}
The three leading Sylvester minors of $M-(7/250)\Id$ are
\begin{equation}
 \frac9{125},\qquad \frac{1211}{250000},\qquad
 \frac{89}{31250000},
 \label{eq:sylvester}
\end{equation}
so $M>(7/250)\Id$.  Parseval therefore gives
\begin{equation}
 \cF_{\rm quad}(\Psi)\ge\frac\nu2\norm{\Psi}_2^2,qquad
 \nu=\mu_{\rm sh}+\frac7{250}.
 \label{eq:nu}
\end{equation}

The full Class-II matrix is positive definite because, after writing
$P=M_X^2+\eps_X>0$, its determinant factor is
\begin{equation}
 c_{JJ}c_{KK}-c_{JK}^2=\frac1{50}>0.
 \label{eq:classii-det}
\end{equation}
For the local polynomial, direct coefficient matching gives
\begin{equation}
 \frac\nu2\rho+\frac\lambda4\rho^2+\frac\gamma6\rho^3
 =g\rho+\frac\gamma6\rho(\rho-\rho_*)^2,
 \quad \rho_*=-\frac{3\lambda}{4\gamma}=\frac{43}{216},
 \quad g=\frac\nu2-\frac{3\lambda^2}{32\gamma}.
 \label{eq:neutral-completion}
\end{equation}
The exact rational evaluation of $g$ is the first value in \eqref{eq:g-kappa}.
After integration and addition of the nonnegative Class-II term,
\eqref{eq:neutral-bounds}'s first inequality follows.  Equality forces
$\rho=0$ almost everywhere, so the zero field is the unique global minimizer.

\subsection{Radial derivative and absence of nonzero equilibria}

The energy lower bound alone does not rule out nonzero saddles.  We therefore
scale $\Psi\mapsto t\Psi$ and put $y=t^2$.  At a fixed field, let $j$ denote
the unscaled current and let $\ell(y)=j-q(y)\nabla\rho$ with
\begin{equation}
 q(y)=\frac{y m}{y\rho+\eps_\rho},qquad
 \theta=\frac{\eps_\rho}{y\rho+\eps_\rho}\in[0,1].
 \label{eq:theta}
\end{equation}
Direct differentiation gives
\begin{equation}
 \frac{\dd E_{\rm II}}{\dd y}
 =y\begin{pmatrix}j&\ell\end{pmatrix}
 R_\theta\begin{pmatrix}j\\\ell\end{pmatrix},qquad
 R_\theta=\begin{pmatrix}
 a-\theta b&b+\theta(b-c)/2\\
 b+\theta(b-c)/2&c(1+\theta)
 \end{pmatrix}.
 \label{eq:Rtheta}
\end{equation}
For the pinned coefficients,
\begin{equation}
 a-\theta b\ge\frac{21}{2000P}>0,qquad
 \det R_\theta=\frac{9(-81\theta^2+128\theta+128)}{10240000P^2}>0
 \quad(0\le\theta\le1).
 \label{eq:Rtheta-positive}
\end{equation}
The numerator is concave and positive at both endpoints, so the Class-II
energy is nondecreasing along every amplitude ray.  Differentiating the
remaining terms in the same ray gives
\begin{equation}
 \ip{D\cF(\Psi)}{\Psi}_{\R}
 \ge\left(\nu-\frac{\lambda^2}{4\gamma}\right)\norm{\Psi}_2^2
 =\kappa\norm{\Psi}_2^2.
 \label{eq:radial-bound}
\end{equation}
This is the second inequality in \eqref{eq:neutral-bounds}.  A nonzero
critical point would make its left-hand side vanish, contradicting
\eqref{eq:radial-bound}.  Finally,
\begin{equation}
 \frac12\frac{\dd}{\dd t}\norm{\Psi(t)}_2^2
 =-\ip{D\cF(\Psi(t))}{\Psi(t)}_{\R}
 \le-\kappa\norm{\Psi(t)}_2^2
 \label{eq:decay-diff}
\end{equation}
proves \eqref{eq:decay}.

\begin{remark}[meaning of the neutral result]
The result rejects nonzero equilibria only for the declared unconstrained,
neutral, hash-pinned functional and its canonical gradient flow.  It does not
reject every retuned functional, fixed-charge problem, compact target, or
conserved evolution.  In particular, it is a candidate-functional rejection
boundary, not a proof that a physical vacuum is empty.
\end{remark}

\section{Exact imposed-ensemble transition}
\label{sec:ensemble-proof}

\subsection{The finite-torus spectral bottom}

The quadratic operator has the form
\begin{equation}
 \mathcal L=(r+Z(-\Delta)+Y(-\Delta)^2)\Id+M,
 \quad M=\begin{pmatrix}
 1/10&-1/20&-1/20\\-1/20&13/100&-1/20\\-1/20&-1/20&17/100
 \end{pmatrix}.
 \label{eq:ensemble-L}
\end{equation}
The characteristic polynomial of $M$ is
\begin{equation}
 p(t)=t^3-\frac25t^2+\frac{223}{5000}t-\frac3{3125}.
 \label{eq:charpoly}
\end{equation}
An exact Sturm computation isolates its smallest root $m_*$ by
\begin{equation}
 \frac{2811617233511}{10^{14}}<m_*<
 \frac{2811617233512}{10^{14}}.
 \label{eq:mstar}
\end{equation}
For $n\in\mathbb Z^3$, $|k_n|^2=\pi^2|n|^2/64$.  The rational enclosure of
$\pi$ used by the audit proves
\begin{equation}
 \frac52<\frac{-Z/(2Y)}{\pi^2/64}<\frac72.
 \label{eq:radial-location}
\end{equation}
Because $|n|^2$ is integral and $3=1^2+1^2+1^2$, the radial minimum is the
shell $|n|^2=3$, containing the eight body-diagonal indices
$n=(\pm1,\pm1,\pm1)$.  Thus
\begin{equation}
 \lambda_0=D(3\pi^2/64)+m_*,qquad
 0.28811617234458494031<\lambda_0<0.28811617234459494032.
 \label{eq:lambda0}
\end{equation}
This is a finite-torus commensurability statement.  It is not a BCC theorem
and it does not remove the shell degeneracy in an infinite-volume limit.

\subsection{Nonnegative decomposition and saturation}

Set $Q=\norm{\Psi}_2^2/2$.  Coefficient matching for
$f(\rho)=\lambda\rho^2/4+\gamma\rho^3/6$ gives
\begin{equation}
 \frac{\lambda_0}{2}\rho+f(\rho)
 =\frac{\mu_t}{2}\rho+\frac{\gamma}{6}\rho(\rho-\rho_*)^2,
 \quad \rho_*=-\frac{3\lambda}{4\gamma},
 \quad c_* =\frac{3\lambda^2}{16\gamma}.
 \label{eq:polynomial-completion}
\end{equation}
Adding the spectral and Class-II remainders yields \eqref{eq:ensemble-decomp}.

Let $u_*$ be a unit eigenvector of $M$ for $m_*$ and select any one of the
eight indices on the shell.  The plane wave
\begin{equation}
 \Psi_*(x)=\sqrt{\rho_*}\,u_*\exp\!\left(\frac{2\pi i}{16}n\cdot x\right)
 \label{eq:plane-wave}
\end{equation}
has constant $\rho$ and constant internal bilinears.  Hence $J_A=K_A=0$,
the spectral remainder vanishes, and
\begin{equation}
 Q_*=\frac{16^3\rho_*}{2}=\frac{11008}{27}.
 \label{eq:Qstar-again}
\end{equation}
All terms in \eqref{eq:ensemble-decomp} are therefore saturated.

\subsection{Fixed charge and grand potential}

For a prescribed charge $Q$, let $\bar\rho=2Q/16^3$.  The exact Bregman
identity is
\begin{equation}
 f(\rho)-f(\bar\rho)-f'(\bar\rho)(\rho-\bar\rho)
 =\frac{(\rho-\bar\rho)^2}{12}
   \bigl(4\gamma\bar\rho+2\gamma\rho+3\lambda\bigr).
 \label{eq:bregman}
\end{equation}
If $\bar\rho\ge\rho_*$, the bracket is nonnegative.  On the affine constraint
$\int\rho=2Q$, the linear term integrates to zero.  The spectral remainder,
the Class-II energy, and \eqref{eq:bregman} consequently prove constrained
global minimality of a constant-density ground-shell plane wave for all
$2Q/16^3\ge\rho_*$.  The constrained problem below $Q_*$ is not decided here.

Subtracting $\mu Q$ from \eqref{eq:ensemble-decomp} gives
\begin{equation}
 \cO_\mu(\Psi)=(\mu_t-\mu)Q(\Psi)
 +\frac12\ip{\Psi}{(L-\lambda_0)\Psi}_{\R}
 +\cF_{\rm II}(\Psi)
 +\frac\gamma6\int\rho(\rho-\rho_*)^2\dd x.
 \label{eq:grand-decomp}
\end{equation}
For $\mu<\mu_t$, every nonzero field has positive grand potential.  At
$\mu=\mu_t$, zero and the saturated states coexist.  For $\mu>\mu_t$, the
plane wave has negative grand potential and sextic coercivity supplies a
global minimizer, which is necessarily nonzero.

The zero-field Hessian loses positivity at $\mu_{\rm lin}=\lambda_0$.  Along
the ground-plane-wave amplitude family, stationary densities solve
\begin{equation}
 \gamma\rho^2+\lambda\rho+(\lambda_0-\mu)=0.
 \label{eq:amplitude-quadratic}
\end{equation}
The amplitude saddle-node is $\mu_{\rm sn}=\lambda_0-\lambda^2/(4\gamma)$,
while coexistence is $\mu_t=\lambda_0-3\lambda^2/(16\gamma)$.  With the
pinned coefficients this is exactly \eqref{eq:first-order-ordering}; the
global charge jumps from zero to $Q_*$ while zero remains a strict local
minimum.  This is the precise sense in which the imposed ensemble transition
is first order.

\begin{remark}[why R-157 and R-158 do not conflict]
At the nonzero state \eqref{eq:plane-wave}, the constrained Euler--Lagrange
equation is $D\cF(\Psi_*)=\mu_t\Psi_*$, not $D\cF(\Psi_*)=0$.  R-157 concerns
the latter equation in the unconstrained linear space; R-158 changes the
variational problem.  Moreover,
$\cF(\Psi_*)=\mu_tQ_*>0=\cF(0)$, so the ordered state does not beat the
original neutral reference.
\end{remark}

\section{Relation to prior work and residual contribution}
\label{sec:prior-art}

The functional-analytic ingredients used here are standard: sectorial
fourth-order semigroups, Sobolev product estimates, Galerkin compactness,
weakly singular Gronwall inequalities, and polynomial coercivity.  A
Swift--Hohenberg analysis by Giorgini proves well-posedness and long-time
attractor properties for scalar slow/fast equations, but does not contain the
present three-component regularized Class-II map, its exact radial sign, or
the imposed-ensemble completion.  Thus the well-posedness theorem alone is
not claimed as a new general method.

Quantum anharmonic-crystal phase-transition work, such as
Kargol--Kondratiev--Kozitsky, develops Euclidean Gibbs states, reflection
arguments, and infrared estimates for specified lattice models.  The present
paper does not import a quantum phase theorem.  It instead treats a classical
fixed-torus variational functional and proves the model-specific sign and
spectral identities needed for its two ensembles.  The residual contribution
is therefore the exact combination of (a) the full derivative Class-II
well-posedness argument, (b) the global neutral rejection, and (c) the finite
shell ensemble classification, all with a common pinned functional and an
explicit scope firewall.

The claim is not a world-first assertion.  A complete novelty determination
requires a specialist literature search beyond the bounded crosswalk recorded
in \texttt{literature-crosswalk.md} and independent mathematical review of the
model-specific estimates.

\section{Verification and reproducibility}
\label{sec:verification}

The paper cites only the registered A2 claim and its result records:
\begin{center}
\small
\begin{tabular}{lll}
\toprule
component & independent executable evidence & recorded result\\
\midrule
full evolution & four audits & 20/20, 14/14, 12/12, 15/15 PASS\\
R-157 neutral sign & primary and non-importing lanes & 26/26, 24/24 PASS\\
R-158 ensemble sign & primary and standard-library lanes & 35/35, 24/24 PASS\\
R-472 exact core & Lean sidecar and hostile checks & 30/30, 22/22, 12/12, 22/22 PASS\\
\bottomrule
\end{tabular}
\end{center}

The canonical commands are listed in \texttt{verification/README.md}.  The
tests recompute rational constants, matrix signs, spectral intervals,
normalizations, and hostile mutations.  They do not prove the analytic
semigroup, compactness, chain-rule, or external-literature steps; those remain
the proof text's responsibility.

The existing A2 T6 reproduction bundle is recorded in the claim card and its
bundle directory.  A dedicated capstone bundle containing the integrated
R-157/R-158 referee note and all transitive scripts is a required
post-review action.  Under the repository policy it must be built only after
operator confirmation of the integrated referee package.  This draft does not
claim that confirmation or a release-gated submission.

\section{Limitations and falsifiers}
\label{sec:limitations}

The theorem's domain is the fixed side-$16$ torus, six-real-component $H^2$
field space, pinned coefficients, positive density floor, and $\eta_{\rm
shell}=0$.  The following are outside its scope:
\begin{itemize}
\item infinite-volume or continuum removal, counterterm flow, and regulator
      independence;
\item a real-time infinite-volume dynamics, algebraic KMS states, ground
      states, spectral gaps, extremality, purity, or clustering;
\item a conserved physical charge, a reservoir, or a derivation of $\mu$;
\item gauge completion, gauge-invariant density modulation, BCC or unique
      morphology;
\item the historical non-variational backend, signed stochastic composites,
      alternative parameters, and compact target spaces;
\item any physical-vacuum, gravity, cosmology, or Sector-A closure claim.
\end{itemize}

The analytic theorem is falsified by a counterexample to global $H^2$
well-posedness, uniqueness, continuous dependence, the exact energy identity,
or positive-time smoothing in the declared domain.  R-157 is falsified by an
exact nonzero critical point or a field violating \eqref{eq:neutral-bounds}.
R-158 is falsified by failure of the exact spectral shell, the decomposition,
the Bregman sign, plane-wave saturation, or the ordering
\eqref{eq:first-order-ordering}.  A failure of the named P1 identification
hypothesis blocks transfer to the TECT production interpretation but does not
alter the mathematics of the explicitly defined functional.

\section{Conclusion}

The pinned functional has a sharp two-sided mathematical verdict.  Its neutral
unconstrained dynamics are globally well posed and converge exponentially to
zero; no nonzero equilibrium survives the exact radial sign.  A nonzero shell
state appears only after the variational problem is changed by an imposed
charge or chemical potential, where an exact first-order coexistence mechanism
can be proved on the finite torus.  This separation is the principal result:
the ensemble mechanism is mathematically real, while its physical provenance
and interpretation remain open.  The manuscript is therefore a candidate
independent mathematical paper, currently at draft status pending the
literature crosswalk, proof audit, operator confirmation, capstone bundle, and
release checks specified in the accompanying files.

\begin{thebibliography}{9}
\bibitem{Pazy}
A. Pazy, \emph{Semigroups of Linear Operators and Applications to Partial
Differential Equations}, Springer, 1983.

\bibitem{Amann}
H. Amann, \emph{Linear and Quasilinear Parabolic Problems}, Vol. I,
Birkhäuser, 1995.

\bibitem{Giorgini}
A. Giorgini, ``On the Swift--Hohenberg equation with slow and fast dynamics:
well-posedness and long-time behavior,'' \emph{Communications on Pure and
Applied Analysis} 15 (2016), 219--241,
\href{https://doi.org/10.3934/cpaa.2016.15.219}{doi:10.3934/cpaa.2016.15.219}.

\bibitem{KKK}
J. Kargol, Y. Kondratiev, and Y. Kozitsky, ``Phase transitions and quantum
stabilization in quantum anharmonic crystals,'' arXiv:0710.2303,
\href{https://arxiv.org/abs/0710.2303}{arXiv:0710.2303}.

\bibitem{AubinLions}
J.-L. Lions, \emph{Quelques méthodes de résolution des problèmes aux limites
non linéaires}, Dunod, 1969.

\bibitem{Temam}
R. Temam, \emph{Infinite-Dimensional Dynamical Systems in Mechanics and
Physics}, 2nd ed., Springer, 1997.
\end{thebibliography}

\end{document}
